\chapter{Example Tcl Client}
	The Tcl client editor.tcl was written by Michael Gebhardt in 47 lines to demonstrate
	 the Tcl interface.
	``cbwish -f editor.tcl'' starts the short demo, if the environment variables are set in the right way. 
	For a quick start you can go to the directory of the example client 
	\begin{verbatim}
		$CB_HOME/examples/Clients/TclClient/ 
	\end{verbatim}
	and start the predefined script ``../../../bin/CBwish -f editor.tcl''. \\
	The first you will see is a standard connection dialog. 
	Please enter the hostname and the portnumber of a running server.
	If the connection was succesful, the dialog will disappear 
	 and two new windows will be opened.
	The first is an editor which can tell its contents to the CBserver 
	 or asks the contents if it calls an instance of ``QueryClass''.
	The second window will display the answers of any query 
	 or the errormessage of the server if a problem occurs.\\
	As a special feature you could double press the left mousebutton on any label
	 in the first editor.
	 If the word was not a label of an object an errormessage will be displayed.
	  Else the contents of the editor will be replaced by the telos frame of this object. 
	 This enables to browse through the frames.
	 
	

\begin{verbatim}
01: #!../bin/cbwish -f
\end{verbatim}
	Line 2 hides the toplevel window.
	Line 3 calls a standard dialog to create a connection to CBserver. 
	The return value is the name of the cb\_server object, 
	 which is stored into the global variable cb.
	If no connection was made, the application will be closed.
\begin{verbatim}
02: wm iconify .
03: set cb [cb_ConnectDlg ""]
04: if { $cb == ""} exit
\end{verbatim}
	The toplevel window ``.'' should contain an editor 
	 and a frame with three buttons for the
	 functions tell, ask and quit.
	A second toplevel should contain a single editor
	 to expose the answers and errormessages.
\begin{verbatim}
05: frame .fButtons 
06: toplevel .fOut
07: text .text -yscrollcommand ".scroll set"
08: text .fOut.text -yscrollcommand ".fOut.scroll set"
09: scrollbar .scroll -command ".text yview"
10: scrollbar .fOut.scroll -command ".fOut.text yview"
11: button .fButtons.bTell -text "tell" -command tell -width 12
12: button .fButtons.bAsk -text "ask" -command ask -width 12
13: button .fButtons.bQuit -text "quit" -command "$cb disconnect; exit" -width 12
\end{verbatim}
	The pack commands organizes the layout
	 of the above constructed widgets by the 
	 geometry manager of Tcl/Tk:
\begin{verbatim}
14: pack .fButtons  -fill x
15: pack .fButtons.bTell .fButtons.bAsk .fButtons.bQuit -side left -pady 10 -padx 5
16: pack .scroll .text .fOut.scroll .fOut.text -side left -fill y
\end{verbatim}
	Describes the titles of both windows and make them visible:
\begin{verbatim}
17: wm title . "ConceptBase Editor 47"
18: wm title .fOut "ConceptBase Editor 47: answers"
19: wm deiconify .
20:
\end{verbatim}
	The ``tell'' procedure will be called by the button constructed in line 11 without
	 any parameters. This function takes the whole contents of the editor and tells it
	 to the CBserver.\\
	These server access is used in line 23. The cb\_server object is expressed by
	the global variable ``cb'' which was set in line 3. 
	Note that the declaration in line 22 is important, 
	 else Tcl tries to evaluate the variable as local.
	\footnote{
		If you create the cb\_server object and connect it yourself 
	 	 without using the standard dialog 
	 	then you could replace the variable by a constant name.
	}
	The text interpreted as telos frame reaches form line 1 position 0 (``1.0'') 
	 to the ``end'' of the editor text.
	If this was wrong an error occurs, 
	 this means the return value will be ``0'' instead of the usual ``1'' 
	 and the answer window will be filled with an errormessage.
	This will be asked by the method ``lastErrormessage'' in line 25.
	At least the answer should remain visible, what ensures line 26.
\begin{verbatim}
21: proc tell { } {
22:     global cb
23:     if [$cb tell [.text get 1.0 end]] {
24: 	    .fOut.text insert end "[.text get 1.0  end]\n"
25:     } else { .fOut.text insert end [$cb lastErrormessage] }   
26:     .fOut.text yview -pickplace end
27: }
28: 
\end{verbatim}
	The ``ask'' procedure will be called by the button constructed in line 12 without
	 any parameters too. This function takes the whole contents of the editor and asks it
	 from the CBserver like the tell procedure explained above.\\
	The difference is the ask method used in line 31. ``answer'' the name of the
	 answervariable, which will be created and assigned by the ask method.
	The contents of this variable will be inserted into the answer window and 
	 additionally echoed on the shell in line 33.
	The last argument ``TELOS'' of the ask method is optional, 
	 but the default format creates an cb\_ptelos objects, 
	 which is useful in normal applications, 
	 but makes no sense in this case.	
\begin{verbatim}
29: proc ask { } {
30:     global cb
31:     if [$cb ask answer [.text get 1.0 end] TELOS] {
32: 	    .fOut.text insert end "$answer\n"
33: 	    puts $answer
34:     } else { .fOut.text insert end [$cb lastErrormessage] }
35:     .fOut.text yview -pickplace end
36: }
37: 
\end{verbatim}
	This load procedure will be bound to the double press 
	 of the left mouse button in line 47.
	The argument should label an object which will be loaded from the database.
	If it doesn't exists the usual errormessage will be displayed. \\
	If you want to ask own Queries then you must prepare them. 
	This means you must tell the query as instance of ``QueryClass''.
	The interface can only handle such told queries, 
	 because the case of the ask procedure above, 
	 that the user formulates the query dynamically, 
	 is hardly to handle.
	For example if you expect an answer of the type FRAME, 
	 but the user asks ``exists[x/objname]'' then an error will occur, 
	 because the format was wrong.
\begin{verbatim}
38: proc load { obj } {
39:     global cb
40:     if { $obj != "" } then {
41:         if [$cb ask answer "get_object\[$obj/objname\]" TELOS] then { 
42: 	        .text delete 1.0 end
43: 	        .text insert end $answer
44: 	    } else { .fOut.text insert end [$cb lastErrormessage] }
45:     }
46: }
\end{verbatim}
	The following statement binds the ``load'' procedure 
	 to the double press of the left mouse button in the editor.
	The argument will be evaluated as the word,
	  wich was clicked by the mouse. \\
	This enables browsing through the telos frames by double pressing.
\begin{verbatim}
47: bind .text <Double-ButtonPress-1> {load [.text get "@%x,%y wordstart" "@%x,%y wordend"]}
\end{verbatim}
